\begin{titlepage}
\thispagestyle{empty}
\vspace*{30mm}
{\sffamily\bfseries\fontsize{34}{40}\selectfont\color{ink}Рой над единой средой\par}
\vspace{10pt}
{\sffamily\fontsize{16}{20}\selectfont\color{midgray}Руководство инженера по протоколу\\распределённой агентной разработки\par}
\vspace{22mm}
{\normalsize\color{ink}
Что это за проект, какие проблемы он решает,\\
из каких примитивов собирается\\
и с чего начинается его реализация\par}
\vfill
{\normalsize\color{ink}
По материалам работы\\
\textit{Разработка программного обеспечения роем автономных агентов над единой средой исполнения}\\
Алексей Григорьев, Иван Опарин, Relux Works, третья редакция, версия 3.3\par}
\vspace{8mm}
{\sffamily\footnotesize\color{midgray}Черновик 0.4, сентябрь 2026 \quad github.com/relux-works\par}
\end{titlepage}

\chapter*{Как читать эту книгу}
\addcontentsline{toc}{chapter}{Как читать эту книгу}

Эта книга написана для инженера, который присоединяется к проекту. Она предполагает десять и больше лет опыта в программировании, хорошее понимание того, как устроены современные агентские системы поверх больших языковых моделей, и не предполагает опыта с распределёнными базами данных. Всё, что нужно знать о репликации, кворумах, арендах и деревьях Меркла, объясняется по ходу в отдельных врезках.

Книга состоит из трёх частей. Первая отвечает на вопрос зачем: какую проблему мы решаем и почему существующие инструменты её не решают. Вторая описывает, как устроен протокол: среда исполнения, журнал, полномочия, оркестратор, цикл задачи, рабочие ноды, модель угроз. Третья описывает, как это построить: из каких существующих компонентов собирается система, что нужно написать заново, как выглядит план на двенадцать недель и какие вопросы остаются открытыми.

Источник для книги это третья редакция работы (версия 3.3), на которой основан протокол, два независимых ревью, которые её сформировали, и обсуждения, которые к ней привели. Там, где книга уходит в детали реализации глубже, чем работа, это отмечено. Там, где решение ещё не принято, это сказано прямо.

Версия 0.4 выравнивает книгу с системой Apiary на 29 сентября 2026 года. Протокол книги есть распределённый слой этой системы, и большинство его понятий теперь принадлежат спецификациям модулей; где книга и спецификация расходятся, права спецификация. Карта разделов работы на их нынешних владельцев дана во введении работы. Главы 10--12 заменены роадмапом экосистемы и репозиториями модулей и сохранены как история, открытые вопросы главы 13 указывают, где вопрос живёт сейчас, а карта имён для людей вынесена в глоссарий.

В тексте используются пять видов врезок.

\begin{likbez}
Объяснение понятия из распределённых систем или криптографии, которое нужно для понимания следующего раздела. Если вы знакомы с темой, врезку можно пропустить.
\end{likbez}

\begin{whybox}
Обоснование архитектурного решения: какие альтернативы были и почему выбрана именно эта.
\end{whybox}

\begin{invbox}
Формулировка одного из десяти инвариантов протокола. Спецификации нижнего уровня ссылаются на инварианты по номерам I1 и далее.
\end{invbox}

\begin{warnbox}
Место, где легко ошибиться при реализации, или свойство, которое протокол не обещает.
\end{warnbox}

\begin{exbox}
Конкретный пример: сценарий, последовательность записей, набросок структуры данных.
\end{exbox}

Терминология выровнена с работой и с существующими спецификациями компонентов. Одно слово, эпоха, в разных компонентах означает разные вещи; в этой книге оно всегда означает срок аренды оркестратора, а остальные значения оговорены там, где встречаются. Глоссарий вынесен в приложение.
